\documentclass[11pt]{article}
\usepackage{mathblatt}
% gesetzt von werkzeuge/setzer.py; Lösungen nach bau/bauregeln.md „Lösungen“.
\newcommand{\kennung}{4CR}
\begin{document}
\pfheftstil
\fancyfoot[C]{{\footnotesize\color{mbgrau}\kennung\hfill\thepage}}
\pfheftkopf{Kenngrößen}{Klasse 7 \textperiodcentered{} Lösungen}

\begin{pfloesung}{}
\lz{1b)}{Min.\ $15$, Max.\ $45$, Spannweite $30$\,min}{Minimum $15$ min, Maximum $45$ min; Spannweite $= 45 - 15 = 30$ min}{}
\end{pfloesung}
\begin{pfloesung}{}
\lz{2.}{Spannweite $0{,}20$\,€, Modalwert $0{,}50$\,€}{Spannweite $= 0{,}60 - 0{,}40 = 0{,}20$ €; Modalwert $= 0{,}50$ € (zweimal)}{}
\end{pfloesung}
\begin{pfloesung}{}
\lz{3.}{Super ($0{,}09$\,€ gegen $0{,}06$\,€)}{Super: Spannweite $= 1{,}85 - 1{,}76 = 0{,}09$ €; Diesel: $1{,}72 - 1{,}66 = 0{,}06$ €. Der Preis für Super schwankte stärker.}{}
\end{pfloesung}
\begin{pfloesung}{}
\lz{4.}{Aussage 2; richtig: Modalwert $9$ Jahre}{– Modalwert $= 9$ Jahre (dreimal), $13$ ist das Maximum; Aussage 1 stimmt: $13 - 9 = 4$}{}
\end{pfloesung}
\begin{pfloesung}{}
\lz{5b)}{$6$ Treffer}{geordnet $3, 4, 6, 7, 9$: Median $= 6$ Treffer}{}
\end{pfloesung}
\begin{pfloesung}{}
\lz{6.}{$13$\,cm}{geordnet $9, 11, 12, 14, 16, 20$: Median $= (12 + 14) : 2 = 13$ cm}{}
\end{pfloesung}
\begin{pfloesung}{}
\lz{7.}{$19{,}5$\,°C}{geordnet $17, 18, 19, 20, 23, 25$: Median $= (19 + 20) : 2 = 19{,}5$ °C}{}
\end{pfloesung}
\begin{pfloesung}{}
\lz{8.}{Ali ($14{,}5$ gegen $14$)}{Tim: geordnet $9, 10, 12, 14, 15, 18, 21$, Median $= 14$; Ali: geordnet $8, 11, 13, 16, 17, 22$, Median $= (13 + 16) : 2 = 14{,}5$. Ali liegt vorn.}{}
\end{pfloesung}
\begin{pfloesung}{}
\lz{9b)}{$2$ Tore pro Spiel}{Summe $= 10$; $10 : 5 = 2$ Tore pro Spiel}{}
\end{pfloesung}
\begin{pfloesung}{}
\lz{10.}{$4{,}11$\,m}{Summe $= 16{,}44$ m; $16{,}44 : 4 = 4{,}11$ m}{}
\end{pfloesung}
\begin{pfloesung}{}
\lz{11.}{$\approx 24\,258$ Zuschauer}{$412\,390 : 17 = 24\,258{,}2\ldots \approx 24\,258$ Zuschauer pro Spiel}{}
\end{pfloesung}
\begin{pfloesung}{}
\lz{12.}{Nein, nur $\approx 193$ Besucher}{Nein: Summe $= 150 + 100 + 150 + 150 + 200 + 300 + 300 = 1\,350$; $1\,350 : 7 \approx 192{,}9$, also weniger als $200$.}{}
\end{pfloesung}
\begin{pfloesung}{}
\lz{13b)}{$10$}{Summe $= 8 \cdot 3 = 24$; dritte Zahl $= 24 - 5 - 9 = 10$}{}
\end{pfloesung}
\begin{pfloesung}{}
\lz{14.}{$64$\,s}{Summe $= 62 \cdot 4 = 248$ s; bekannt $60 + 65 + 59 = 184$ s; vierte Runde $= 248 - 184 = 64$ s}{}
\end{pfloesung}
\begin{pfloesung}{}
\lz{15.}{$11$ Jahre, es sinkt}{Summe $= 12 \cdot 5 = 60$; neue Summe $= 60 - 16 = 44$; $44 : 4 = 11$ Jahre; das Mittel sinkt}{}
\end{pfloesung}
\begin{pfloesung}{}
\lz{16.}{Ja, mit mindestens $17$ Punkten}{Ja: Summe für fünf Tests $= 15 \cdot 5 = 75$; bisher $14 + 17 + 12 + 15 = 58$; nötig $75 - 58 = 17$ Punkte, das ist höchstens $20$.}{}
\end{pfloesung}
\begin{pfloesung}{}
\lz{17.}{$20$\,°C}{geordnet $16, 17, 19, 21, 22, 24$: Median $= (19 + 21) : 2 = 20$ °C}{}
\end{pfloesung}
\begin{pfloesung}{}
\lz{18.}{Do: $95$ Besucher}{Summe $= 120 \cdot 6 = 720$; bekannt $95 + 130 + 110 + 140 + 150 = 625$; Donnerstag $= 720 - 625 = 95$ Besucher}{}
\end{pfloesung}
\begin{pfloesung}{}
\lz{19.}{Aussage 2; richtig: Median $40$}{– geordnet $37, 38, 38, 39, 40, 40, 40, 41, 42$, der Median ist der $5.$ Wert: Median $= 40$; Aussage 1 stimmt: $42 - 37 = 5$}{}
\end{pfloesung}
\begin{pfloesung}{}
\lz{20.}{\mbox{a) $46$ Jahre} \mbox{b) Nein, bleibt $46$ Jahre}}{a) Summe $= 506$; $506 : 11 = 46$ Jahre; b) Nein: $506 - 68 - 24 = 414$; $414 : 9 = 46$ Jahre. Die beiden sind zusammen $92 = 2 \cdot 46$ Jahre alt.}{}
\end{pfloesung}
\end{document}
